\section{Deployment Recommendation, Limitations, and the Pre-Registered Live Cohort}\label{sec:discussion}
\subsection{Recommendations from the simulator, and the protocol-only tier}\label{sec:deploy}
The recommendations below have two tiers: what this paper's simulator
measures, and what remains protocol-only.

\paragraph{Mechanism-level (measured in this paper's simulator).}
Where a service already requires protected verification evidence and a
constrained channel, the serialized-evidence component of
Section~\ref{sec:substrate} applies as before. The
retirement study adds a sharper pattern for context policy: in the
simulator, the policy that first converts history into current
verified bindings, a typed pending-obligation set, and dependency
interfaces passes the acceptance contract at every scale tested,
while policies that rely on lossy summaries of obligations, or that
revert values to beliefs older than the current working window,
measurably fail (Section~\ref{sec:retire-results}). The
conversion is the pattern the simulator supports; whether it is worth
adopting depends on the
cost regime and on the retrieval layer. It buys \RetwoOpsReduction{}
fewer retrieval round trips at equal correctness, and by the closed
form of Proposition~\ref{prop:breakeven} the simulator's gate flips
at a computable round-trip price: given means
$(B_f,B_m,o_f,o_m)$ measured in a deployment's own pilot and that
deployment's price $c$, the formula's threshold $c_{0.8}$ states
where the frozen 20\% gate would flip. In the simulator,
$c_{0.8}\approx1{,}364$\,B pooled (frontier's 8.0\% byte advantage
does not by itself clear the 20\% gate; the frozen 4{,}096\,B
round-trip price does). These are simulator thresholds under the
frozen accounting, not validated deployment guidance: they identify
what a deployment must measure---its own byte costs, operation mix,
and round-trip price---rather than prescribing adoption.
Second, and now stated because the exploratory analyses forced it: a
64-gate bounded-recency window (\texttt{bounded-hier}) is cheaper
than frontier below a break-even price of
\ExpBhBreakEvenPooled{} pooled (2{,}148\,B independent; never in
global) at equal acceptance---but \ExpBhDueExposed{} of its due
obligations and \ExpBhFootprintExposed{} of its footprint reads sit
outside its retention, so its safety delegates to the retrieval
layer. The simulator-conditional reading is therefore a function of
two observable quantities, not one: with reliable, cheap retrieval and
round-trip price below the break-even, the bounded window is the
cheaper arm in the simulator; with expensive or unreliable retrieval,
the conversion heuristic is the arm whose model-internal guarantees
do not depend on the retrieval layer beyond the
fetch itself. Which of these a real deployment should prefer is
exactly what the live cohort measures; the simulator brackets the
possibilities rather than deciding between them. The real-trace replay
(Section~\ref{sec:retire-exploratory}) brackets what this looks like
on non-synthetic change patterns. Between grouped
certificate layouts and flat delta ledgers, select by update density:
the measured-density rule matched the per-cell oracle on every cell of
the frozen cohort (\RetwoDensityRegret{} regret); dense updates favor
the tree, sparse deltas favor the flat ledger, and the density signal
is observable on a training half of any real stream. The
demand-driven adaptive variant is not promoted: its pre-registered
gate failed in the simulator
($\RetwoAdaptiveReduction{}$ versus plain frontier), and adoption
decisions belong to the live cohort, not to the simulator alone.

\paragraph{Live-agent cohort (pre-registered, not run; highest-priority next experiment).}\label{sec:live-protocol}
The mechanism study predicts, but cannot confirm, that avoided
retrieval operations translate into live savings. Executing this
cohort is the single most informative next step the study permits
itself: it is the only pre-registered instrument that can move
real-world agent relevance from moderate to measured, and every
limitation this section records---BPE token counts, provider billing,
prompt-semantics effects, and live tool-call latencies---is either an
outcome it measures or a risk it prices. The confirmatory
cohort is therefore pre-registered with operations as the primary cost
measure:
\begin{enumerate}[leftmargin=*,itemsep=1pt]
\item \textbf{Workload.} Multi-stage changes on a fixed public issue
benchmark (SWE-bench-class), screened for dependency coupling; the
same frozen independent acceptance contract for all arms; equal total
budgets, sandbox, and protected state; every failed branch in the
numerator.
\item \textbf{Arms.} The arms the simulator separates:
observation masking; summary/reset; verified-frontier retirement; and
the density-selected evidence layout (flat delta ledger versus
grouped tree, chosen per the measured update density of the pilot).
The demand-driven adaptive variant is excluded because its promotion
gate failed in the simulator; including it would re-litigate a
pre-registered decision.
\item \textbf{Primary outcomes.} Tool-call round trips per accepted
project (the pre-registered live-cost measure, from the mechanism
result); cost per accepted project (all provider metering preserved,
unknowns never zero-imputed); acceptance rate; deadline success with
censored timeouts. Pre-specified noninferiority margin on acceptance:
no arm may trail masking by more than 5\% absolute.
\item \textbf{Closure rules.} \emph{Null closure}: if frontier
retirement does not reduce tool-call operations per accepted project
by at least 20\% versus masking at noninferior acceptance, the
live-cohort hypothesis is rejected and retirement is not promoted.
\emph{HACT closure}: if the density-selected layout does not beat
flat-delta evidence by at least 5\% on measured evidence bytes at
equal correctness, the tree is not deployed.
\end{enumerate}
Two explicit closure rules prevent indefinite iteration; either rule
terminates the program with a written result. This protocol is frozen
here, before any live data collection.

The live layer also prices risks the mechanism study cannot express
and must not silently assume away. A hosted model does not read
idealized state capsules: structured JSON or capsule formats can
induce attention degradation, prompt distraction, or formatting drift,
so the cohort scores each arm's full delivered prompts, not just their
sizes. Real summarization does not follow a synthetic four-episode
boundary, so the summary/reset arm keeps its own provider-side
implementation with the acceptance contract as the only fixed point.
And live round trips are not a static additive constant: network
jitter, rate limiting, and variable generation latency replace the
frozen $4{,}096$\,B accounting constant with the cohort's measured
per-tool-call durations, which is why tool-call counts, not byte
totals, remain the pre-registered primary outcome.

\paragraph{What remains unmeasured.}
No hosted model was called for any number in this paper; BPE
tokenization, physical-WAN behavior, and provider billing remain
unknown and are never zero-imputed. Two tokenization caveats deserve
emphasis because the paper's savings are reported in bytes. First,
both requested BPE encodings (\texttt{cl100k\_base} and
\texttt{o200k\_base}) remain \UNKNOWN{}
(Section~\ref{sec:unmeasured}); byte savings therefore do not
necessarily convert to proportional BPE token savings, and no
byte-to-token division is asserted anywhere. Second, even a measured
conversion would be model-specific: tokenization is not a
length-preserving map, so the live cohort records actual tokenizer
output per arm rather than converting the byte result. The mechanism cohort's cost model
is a frozen accounting of a frozen model of information flow; real
payloads are not lognormal, and real summaries may retain obligations
better or worse than the seeded boundary survival. The round-trip
regime's $4{,}096$\,B constant is an accounting choice bracketed by a
sensitivity range, not a measured latency. The tool-call prediction
is the only bridge from the mechanism result to live cost, and it is
pre-registered as such rather than assumed.

\paragraph{Trust and scope.}
Future Code remains a single-host supervisor. The guard and TLS
utility are not a distributed consensus protocol or fault-tolerant
deployment. A passing declared suite cannot establish absence of
arbitrary defects. Authorization requires an uncompromised checker,
correct protected reference data, immutable candidate identity, and
local monotone state; Theorem~\ref{thm:soundness} stands as before.
TLS authenticates the configured server, not the checker's
truthfulness. The receiver is a bounded measurement utility behind an
operator-controlled firewall, not a production ingestion service.
